Rescues occur in 27 of the 49 failure-containing tasks. BreadSelection and PrepareCoffee contribute 13 of the 56 rescues (23.21\%); the five largest rescue counts contribute 24 (42.86\%). All task rows, including zero-rescue tasks, are retained in the supplement. Equal weighting of each failure-containing task gives a 7.83\% mean rescue fraction, compared with the failure-weighted 5.58\%. These answer different questions: the latter gives more weight to tasks with more historical failures.

The most difficult original tasks remain difficult after assistance. The three tasks with zero historical successes yield only one rescue among 150 failures (one among 135 eligible pairs). The five lowest historical task scores together yield one rescue among 237 failures (one among 219 eligible pairs). Thus the positive outcomes do not remove the dominant hard-task bottleneck (Fig.~\ref{fig:tasks}). These subsets are retrospective descriptions, not independently chosen test cohorts.

Resampling whole tasks within each original split (30,000 draws, fixed seed) gives a 2.5th--97.5th percentile range of 3.59--7.95\% for the failure-weighted confirmed fraction. This is an exploratory task-mixture sensitivity range, not a confidence interval for the fixed cohort or a population treatment effect. Shared assistance batches create an additional dependence not modeled by this resampling; we therefore do not use it for a significance claim.
